\chapter{Analysis of Variance}
\label{ch:anova}

\noindent\textsc{When you compare two groups,} you use a $t$-test. When you compare three or more, you don't run a series of $t$-tests. Each comparison has its own probability of a false positive, and they accumulate. Run twenty independent tests at $\alpha = 0.05$, and you'd expect one significant result by chance even when nothing is going on. This is the \emph{multiple comparisons problem}, and the cleanest solution for comparing several means is analysis of variance.

ANOVA partitions the total variation in the data into a part attributable to differences between groups and a part attributable to variation within groups. If the between-groups variation is large relative to the within-groups variation, the groups differ.

\section{The setup}

You have $k$ groups with sample sizes $n_1, n_2, \ldots, n_k$ and a total of $n = \sum n_j$ observations. Let $y_{ij}$ denote the $i$-th observation in group $j$, $\bar y_j$ the mean of group $j$, and $\bar y$ the grand mean of all observations.

The null hypothesis: all groups have the same population mean,
\[
  H_0\colon \mu_1 = \mu_2 = \cdots = \mu_k.
\]
The alternative: at least one group mean differs from the others. Notice the alternative is not specific about which means differ. ANOVA tells you that something is going on; post-hoc tests tell you what.

\section{The variance decomposition}

The total sum of squares splits into between-groups and within-groups parts:
\begin{align}
  \underbrace{\sum_{j=1}^{k}\sum_{i=1}^{n_j}(y_{ij} - \bar y)^2}_{\text{SS}_{\text{tot}}}
  &= \underbrace{\sum_{j=1}^{k} n_j (\bar y_j - \bar y)^2}_{\text{SS}_{\text{between}}}
  + \underbrace{\sum_{j=1}^{k}\sum_{i=1}^{n_j}(y_{ij} - \bar y_j)^2}_{\text{SS}_{\text{within}}}.
  \label{eq:anova-decomp}
\end{align}
The total variation around the grand mean equals the variation of group means around the grand mean (between) plus the variation of observations around their group means (within). This decomposition is exact. It's the same Pythagorean structure as the OLS variance decomposition you saw earlier.

\section{The $F$-statistic}

Convert the sums of squares into mean squares by dividing by their degrees of freedom:
\begin{align}
  \text{MS}_{\text{between}} &= \frac{\text{SS}_{\text{between}}}{k - 1}, \\[3pt]
  \text{MS}_{\text{within}}  &= \frac{\text{SS}_{\text{within}}}{n - k}.
\end{align}
The ANOVA $F$-statistic is the ratio:
\begin{equation}
  F = \frac{\text{MS}_{\text{between}}}{\text{MS}_{\text{within}}}.
  \label{eq:anova-f}
\end{equation}
Under $H_0$, this follows an $F$-distribution with $(k-1, n-k)$ degrees of freedom. Large $F$ means the between-group variance dominates the within-group variance: the groups really do have different means. Small $F$ means the differences between group means are no larger than what within-group noise would produce.

The intuition: if all groups have the same population mean, the only source of variation in the group means is sampling error. The ratio of between-group to within-group variation should be roughly 1. When the ratio is much greater than 1, the group means must reflect real differences.

\begin{figure}[ht]
\centering
\includegraphics[width=\linewidth]{fig_anova.pdf}
\caption{The ANOVA decomposition. Left: between-group variation, the spread of group means (thick horizontal lines) around the grand mean (dashed line). Right: within-group variation, the spread of observations around their own group mean. ANOVA's $F$-statistic is the ratio of these two variances. When between is large relative to within, the groups differ.}
\label{fig:anova}
\end{figure}



\section{Worked example}

Three teaching methods, eight students each. Final exam scores:

\begin{table}[ht]
\centering
\begin{tabular}{lcccccccc|c}
\toprule
Method & \multicolumn{8}{c|}{Scores} & $\bar y_j$ \\
\midrule
A & 78 & 82 & 75 & 80 & 77 & 81 & 79 & 76 & 78.5 \\
B & 85 & 88 & 82 & 90 & 86 & 84 & 87 & 83 & 85.6 \\
C & 70 & 73 & 68 & 75 & 72 & 71 & 69 & 74 & 71.5 \\
\bottomrule
\end{tabular}
\caption{Test scores under three teaching methods.}
\end{table}

Grand mean: $\bar y = (78.5 + 85.6 + 71.5)/3 = 78.5$.

Between-groups SS: $8 \cdot [(78.5 - 78.5)^2 + (85.6 - 78.5)^2 + (71.5 - 78.5)^2] = 8 \cdot [0 + 50.4 + 49.0] = 795.2$.

Within-groups SS: sum of $(y_{ij} - \bar y_j)^2$ across all 24 observations $\approx 130$.

$F = (795.2 / 2) / (130 / 21) = 397.6 / 6.19 \approx 64.2$.

This is enormous. The critical $F$ at $\alpha = 0.05$ for $(2, 21)$ degrees of freedom is about $3.47$. Reject $H_0$ overwhelmingly. The teaching methods produce different mean scores.

But notice what we don't know yet: which methods differ? All three could differ from each other, or two could be similar and one different.

\section{Post-hoc tests}

After a significant ANOVA, you compare specific pairs of means. Doing this naively reintroduces the multiple-comparisons problem, so several procedures adjust:

\paragraph{Tukey's HSD.} The standard. Compares all pairs while controlling the family-wise error rate. The honest significant difference is
\[
  \text{HSD} = q_{\alpha, k, n-k} \cdot \sqrt{\frac{\text{MS}_{\text{within}}}{n_j}},
\]
where $q$ is the studentized range distribution. Two groups differ significantly if their mean difference exceeds the HSD.

\paragraph{Bonferroni correction.} Run all $\binom{k}{2}$ pairwise $t$-tests, but require $p < \alpha / \binom{k}{2}$ for significance. Conservative; loses power as $k$ grows.

\paragraph{False discovery rate (Benjamini-Hochberg).} Controls the expected proportion of false positives among the rejected hypotheses. Less conservative than Bonferroni and increasingly the default in fields where many comparisons are routine.

\section{ANOVA as a regression}

ANOVA is multiple regression in disguise. Code group membership with $k - 1$ dummy variables (leaving one as the baseline):
\[
  y_i = \beta_0 + \beta_1 D_{1i} + \beta_2 D_{2i} + \cdots + \beta_{k-1} D_{(k-1)i} + \varepsilon_i.
\]
The $\beta_0$ is the mean of the baseline group. Each $\beta_j$ is the difference in mean between group $j$ and the baseline. Testing whether all the $\beta_j$ are zero is exactly the ANOVA $F$-test.

This equivalence matters. It means you can do ANOVA inside any regression framework: with covariates (ANCOVA), with multiple factors (two-way ANOVA), with random effects (mixed models). ANOVA isn't a separate technique. It's a special case of the linear model.

\begin{keyinsight}
ANOVA tests whether the means of three or more groups differ.

It works by partitioning total variance into between-group and within-group components. Large between-group variance relative to within-group variance produces a large $F$-statistic, indicating the groups differ.

A significant ANOVA tells you that something is happening; post-hoc tests (Tukey, Bonferroni, FDR) tell you which groups differ. Equivalent to a regression with dummy variables for groups.
\end{keyinsight}

\begin{codelab}
\begin{lstlisting}[language=Rlang]
# Three teaching methods
scores <- c(78, 82, 75, 80, 77, 81, 79, 76,
            85, 88, 82, 90, 86, 84, 87, 83,
            70, 73, 68, 75, 72, 71, 69, 74)
method <- factor(rep(c("A", "B", "C"), each = 8))

# One-way ANOVA
fit <- aov(scores ~ method)
summary(fit)

# Tukey HSD post-hoc
TukeyHSD(fit)

# Equivalent regression
lm(scores ~ method)

# Built-in dataset example
data(PlantGrowth)
boxplot(weight ~ group, data = PlantGrowth)
summary(aov(weight ~ group, data = PlantGrowth))
\end{lstlisting}
\end{codelab}

\section{Two-way ANOVA}

Two factors at once: teaching method (3 levels) and time of day (morning vs.\ afternoon). The model is
\[
  y_{ijk} = \mu + \alpha_i + \beta_j + (\alpha\beta)_{ij} + \varepsilon_{ijk}.
\]
Three sources of variation: main effect of method ($\alpha$), main effect of time ($\beta$), interaction $(\alpha\beta)$. The interaction tests whether the effect of method depends on time of day.

Two-way ANOVA is just a regression with main-effects dummies and interaction terms. Same machinery, more sources of variation.

\section{Practice problems}

\begin{enumerate}
  \item For the teaching-methods example, compute the within-groups SS by hand. Verify the $F$-statistic. Find the $p$-value using \texttt{pf()}.
  \item Show that for $k = 2$ groups, ANOVA's $F$-test and the two-sample $t$-test produce identical $p$-values, with $F = t^2$.
  \item Run ANOVA on the PlantGrowth dataset in R. Which groups differ according to Tukey's HSD?
  \item Why does running multiple pairwise $t$-tests inflate the family-wise error rate? Show with simulation: generate three groups from the same population, run all pairwise tests at $\alpha = 0.05$, repeat many times.
  \item Set up the regression equivalent of a one-way ANOVA with four groups. How many dummy variables? What does each coefficient mean?
  \item The interaction in two-way ANOVA tests whether the effect of one factor depends on the other. Sketch what main effects and interactions look like graphically (interaction plots).
\end{enumerate}

\begin{reflect}
$\triangleright$~Why is ANOVA preferred over running multiple $t$-tests when comparing several groups?

\smallskip
$\triangleright$~A significant ANOVA tells you the groups differ. What does it not tell you, and what additional procedure is needed?

\smallskip
$\triangleright$~How is ANOVA equivalent to regression with dummy variables? Why does this equivalence matter?
\end{reflect}
